\documentclass{article}
\newcommand{\ip}[2]{(#1, #2)}
\begin{document}
Since any number of consecutive spaces are treated
like a single one, the formatting of the input
file makes no difference to
      \LaTeX,
but it makes a difference to you.  When you use
\LaTeX, making your input file as easy to read
as possible will be a great help as you write
your document and when you change it.  This sample
file shows how you can add comments to your own input
file.

Quotation marks like
       ``this''
have to be handled specially, as do quotes within
quotes:
       ``\,`this'
        is what I just
        wrote, not  `that'\,''.

       Gnats, gnus, etc.\ all
       begin with G\@.
You should check the spaces after periods when
reading your output to make sure you haven't
forgotten any special cases.  Generating an
ellipsis
       \ldots\
with the right spacing around the periods requires
a special command.

\LaTeX\ interprets some common characters as
commands, so you must type special commands to
generate them.  These characters include the
following:
       \$ \& \% \# \{ and \}.

In printing, text is usually emphasized with an
       \emph{italic}
type style.

\begin{em}
   A long segment of text can also be emphasized
   in this way.  Text within such a segment can be
   given \emph{additional} emphasis.
\end{em}

It is sometimes necessary to prevent \LaTeX\ from
breaking a line where it might otherwise do so.
This may be at a space, as between the ``Mr.''\ and
``Jones'' in
       ``Mr.~Jones'',
or within a word---especially when the word is a
symbol like
       \mbox{\emph{itemnum}}
that makes little sense when hyphenated across
lines. Macros: \ip{A}{B} and \TeX.
\end{document}
